\section{Exponential smoothing for every partial deck}
\label{n:strong-providers}

We now extract two reusable consequences of an exponential density
cost per tracked card: a pointwise bound after a fixed number of
compositions, and a contraction factor involving the dimension of
the diagram below the first row.  The density input already follows
from Theorem~\ref{fac:main}, since $(l/N)^\eta\le1$; renaming
coordinates covers every order and both positive orientations.
Here we give an independent proof by encoding cycle closures through
the switch bits they force.  The proof applies uniformly from the
empty tuple to the full deck and makes no representation estimate
until the final lower-diagram transfer.

Let $N=2^d$, with $d\ge0$, and let $T_N$ be one sweep of independent
fair switches in any fixed order of the $d$ distinct cube coordinates.
Permutations map input positions to output positions, and products
act from right to left.  The inverse-permutation law is denoted by
$T_N^*$.  Thus
\[
 H_N\in\{T_N^*T_N,T_NT_N^*\}
\]
is a forward sweep followed by its reversed coordinate order, with
fresh coins in both halves; the choice of the first order distinguishes
the two displayed orientations.  On
\[
 \mathcal X_{N,l}=\{(x_1,\ldots,x_l):x_i\in\{0,1\}^d,
                      \ x_i\ne x_j\ (i\ne j)\},
 \qquad 0\le l\le N,
\]
write $H_N(x,y)$ for the transition probability, and $u_{N,l}$ for
uniform measure on this set.  Put $(N)_l=N!/(N-l)!$, with $(N)_0=1$.
The empty tuple has transition probability one.

\begin{theorem}[Strong density bound in both orientations]
\label{n:strong-density}
There are absolute constants $p\in(1,2]$ and $C<\infty$ such that,
for every $N=2^d$, every coordinate order, every $0\le l\le N$,
either choice of $H_N$, and every $x\in\mathcal X_{N,l}$,
\begin{equation}\label{n:strong-injection-density}
 \frac1{(N)_l}\sum_{y\in\mathcal X_{N,l}}
       \big((N)_lH_N(x,y)\big)^p\le e^{Cl}.
\end{equation}
After increasing $C$, the same constants give
\begin{equation}\label{n:strong-density-power-normalization}
 N^{-l}\sum_{y\in\{0,1\}^{dl}}
       \big(N^lH_N(x,y)\big)^p\le e^{Cl},
\end{equation}
where the kernel is extended by zero on tuples with repeated positions.
Both conclusions hold with rows replaced by columns.
\end{theorem}

We prove the injection-density bound first.  The power normalization
will follow by an explicit falling-factorial comparison.  The central
counting estimate permits fixed permutations before and after the
fresh layers; this is what allows it to be used after conditioning on
the other parts of the network.

\subsection{Encoding cycles by omitted switch bits}
\label{n:cycle-encoding}

We need a counting observation on cycles. Let \(m\) be a cube size, \(D\) a fixed set of at most \(k\) positions in it, and \(\pi\) the permutation obtained using any fixed pre-permutation, fair switch layers in \(L\) distinct coordinates, and any fixed post-permutation. If \(Y\) is the number of cycles of \(\pi\) contained entirely in \(D\), there are absolute \(c,\gamma>0\) and \(L_0\) such that
\begin{equation}
\Pr(Y\ge t)\le 2^{-c L t}
\qquad
\big(L\ge L_0,\ t\text{ integer}\ge\max(1,k2^{-\gamma L})\big).
\label{n:encoded-cycle-tail}
\end{equation}
The contact estimate of Lemma~\ref{net:encounters} applies after
relabeling inputs by the fixed pre-permutation.  If the $L$ coordinates
do not fill the whole cube, apply it separately in each $L$-coordinate
block: $\sum_iw_i\log_2w_i\le w\log_2w$ shows that any $w$ selected
paths still share at most $w\log_2w/2$ switches in total.

We prove~\eqref{n:encoded-cycle-tail} with $\gamma=1/1000$ and
$L_0=256$.  The event is empty if $t>k$.  Otherwise, when $Y\ge t$,
at least $t/2$ of its cycles have length at most $2k/t$.  Select
$s=\max(1,\lfloor t/2\rfloor)$ of these cycles; in particular
$s\ge t/3$.  We will encode their routes while omitting the fresh
switch bits forced by closing a cycle.

First we find starts and an exploration order that omit many bits.
Choose each start uniformly in its cycle.  Independently give $D$
a uniform random priority order, and visit the selected cycles in
the order of their starts, completing one cycle before the next.
On a cycle of length $w$, the closing path is uniform among its
$w$ paths.  Its expected contacts with other paths of that cycle
are at most $\log_2w$ by the contact estimate.  At each of the $L$
layers, a contact outside that cycle can reveal at most one bit;
conditional on the starts, the contacted selected cycle precedes
it with probability $1/2$.  Therefore the number $S$ of fresh
switch bits on the closing paths satisfies
\[
 0\le S\le Ls,\qquad
 \E S\ge Ls/2-s\log_2(2k/t)\ge0.4Ls.
\]
The last inequality uses $k/t\le2^{\gamma L}$ and $L\ge256$.
It follows that $\Pr\{S\ge Ls/4\}\ge1/5$.  Thus at least one
fifth of the priority orders admit some choices of starts saving
$Ls/4$ bits.  This is an existence statement for each coin
configuration, before counting configurations under fixed priorities.

Fix a priority order.  A code consists of the unordered set of $s$
starts, the cycle lengths in traversal order, and a finite bit string.
The decoder keeps a table of already known switch bits.  It traces
nonclosing edges using the fixed pre- and post-permutations, reading
a bit only at a previously unknown switch.  On a closing edge,
its input and required endpoint are known.  A butterfly updating
each coordinate at most once has at most one route between those
endpoints; the decoder fills the unknown bits along that route with
the forced values, rejecting inconsistencies.  After completing the
cycles, it reads all remaining unknown bits in a fixed order.
Hence a valid code determines at most one full coin configuration.
If $S\ge Ls/4$, its bit string has length at most $mL/2-Ls/4$.
Counting all finite strings up to that length, including unequal
lengths, gives at most
\[
 \binom ks(2k/t)^s\,2^{mL/2-Ls/4+1}
\]
configurations for these fixed priorities.

Every configuration in the tail event has such a code for at least
one fifth of all priority orders.  Averaging the preceding count
over priorities and dividing by the $2^{mL/2}$ equiprobable coin
configurations bounds the tail probability by
\[
 5\,2^{1+s\{\log_2(6e)+2\gamma L-L/4\}}
 \le 2^{-Lt/64}\qquad(L\ge256).
\]
Here
$\log_2\binom ks+s\log_2(2k/t)
 \le s\{\log_2(6e)+2\gamma L\}$ uses $s\ge t/3$.
This proves~\eqref{n:encoded-cycle-tail}.

\subsection{There-and-back recursion on cube nodes}
Now realize \(H_N\) as there-and-back switches (using the pass order or its reverse first, according to which product is used), all with independent coins. It is a binary tree of cube nodes: each size \(2m\) node uses switches first across one axis, then gives two children (panels of \(m\) positions, indexed by columns of the switch pairs) on which independent child there-and-back permutations act, then again switches on its original axis. Size one leaves do nothing. Call the first layers of switches over the tree its input layers, and second layers its output layers.

Within a node, write \(R(x,y)\) for the density of output placement at \(y\) for \(k\) distinguished cards starting at \(x\), relative to uniform on such placements. Write \(J(x,y)\) for the density in an ideal experiment at the node, replacing the two child permutations by independent uniforms. We use at root or any node, for \(p=1+\theta>1\), the inequality
\begin{equation}
\mathbb E_{\rm unif} R(x,y)^p G(y)
\ \le\
\mathbb E_{\rm true}\Big[G(y)\prod_v J_v(X_v,Y_v)^\theta\Big]
\label{n:ideal-recursion}
\end{equation}
for any nonnegative \(G\). Left expectation is in \(y\); right uses the actual switch experiment with this input \(x\), taking \(J_v\) at this and descendant nonleaf nodes with the placements of the distinguished cards passing through them (local inputs \(X_v\) and outputs \(Y_v\)). Densities for empty placements are 1. To see \eqref{n:ideal-recursion}, in the ideal node experiment take \(x_i,y_i\) the inputs into and outputs of children. Given final \(y\) (when \(J>0\)), the conditional expectation of the product of child true densities \(R_i(x_i,y_i)\) equals \(R(x,y)/J(x,y)\), by the density change to the true experiment from ideal. And \(R\) vanishes when \(J\) does. Jensen thus bounds the left side by
\[
\mathbb E_{\rm ideal} G(y)\,J(x,y)^\theta\,\prod_i R_i(x_i,y_i)^p.
\]
Condition on this node's outer switch layers (input and output), so \(y_i\) in this formula are independent uniforms, and \(y\) is a function of them and the outer choices. Apply \eqref{n:ideal-recursion} inductively at each child (against nonnegative functions of their outputs); this gives exactly the desired inequality. The size one case is immediate.

For \(x,y\) in the size \(2m\) node, let the multigraph on the \(k\) cards have two matchings, given by sharing an input column or sharing an output column, of sizes \(a,b\) respectively. It consists of paths (possibly single vertices) and even cycles (possibly of length two); let \(c_*\) be the number of cycles. Call coloring the cards by the two children allowed if opposite along matching edges. With \(r\) the number in the first child, we have
\begin{equation}
J(x,y)=(2m)_k 2^{a+b-2k}
  \sum_{\text{allowed colors}}\frac{1}{(m)_r(m)_{k-r}}
=2^{c_*}\mathbb E Z(r),\qquad
 Z(r)=\frac{(2m)_k}{2^k(m)_r(m)_{k-r}}.
\label{n:node-colors}
\end{equation}
Indeed realizing a given coloring from input costs \(2^{a-k}\); given correct child output columns the final switches cost \(2^{b-k}\); the child-permutation probabilities between columns are as in the sum, with columns distinct per color by the constraints. There are \(2^{k-a-b+c_*}\) allowed choices, taken equally in the last expectation.

Every allowed coloring is feasible: it puts at most one card of
each color in each of the $m$ input columns and each output column.
Each path or cycle component has two colorings.  Even components
contribute equally to the two colors; each odd path component
contributes one independent sign to their difference.

The following bound holds uniformly:
\begin{equation}
\mathbb E Z(r)\le \exp\big(C k\log(2m)/m\big).
\label{n:color-imbalance}
\end{equation}
To prove it, put $z=2r-k$.  Each odd-vertex path component contributes
one independent uniform sign to $z$; let $q$ be their number.  Clearly
$q\le k$.  Making all odd components favor the same color is feasible,
so $k+q\le2m$ and hence $q\le w:=2m-k$ as well.  At the balanced
count $r_0=\lceil k/2\rceil$, writing $k=2a$ or $2a+1$ gives
\[
 Z(r_0)=\prod_{i=0}^{a-1}\frac{2m-2i-1}{2m-2i}\le1.
\]
The cases $k=0$, $q=0$, or $w=0$ are consequently immediate.

Suppose first that $1\le k\le\epsilon m$, for a small absolute
$\epsilon$.  The factorial ratio
\[
 \frac{Z(r)}{Z(r_0)}
 =\frac{(m-r)!(m-k+r)!}
        {\lfloor w/2\rfloor!\lceil w/2\rceil!}
\]
is unchanged by replacing $z$ by $-z$.  Moving $i$ steps from the
balanced count multiplies successive factors of the form
$1+O((i+1)/m)$.  Thus it is at most $e^{C(z^2+1)/m}$.
Since $z$ is a sum of at most $k$ fair signs,
$\Pr\{z^2\ge v\}\le2e^{-v/(2k)}$, and for $m>2Ck$,
\[
 \E e^{Cz^2/m}
 \le1+\frac{4Ck}{m-2Ck}.
\]
Choosing $\epsilon$ small makes $\E Z(r)\le e^{C'k/m}$,
which implies~\eqref{n:color-imbalance} in this range.

For $k>\epsilon m$ and $w>0$, Stirling's formula gives the
uniform bound
\[
 \frac{Z(r)}{Z(r_0)}
 =\frac{((w-z)/2)!((w+z)/2)!}
        {\lfloor w/2\rfloor!\lceil w/2\rceil!}
 \le e^{O(\log(2m))+wI(z/w)},
\]
where
$I(x)=((1-x)\log(1-x)+(1+x)\log(1+x))/2$
is extended continuously at the endpoints.  The sign sum has point
probabilities at most $e^{-qI(z/q)}$.  Convexity and $I(0)=0$ give
\[
 I(z/w)\le(q/w)I(z/q)\qquad(0<q\le w).
\]
The two exponential rate factors therefore cancel.  Summing over
at most $2m+1$ possible values of $z$ leaves $e^{O(\log(2m))}$,
which is at most $e^{Ck\log(2m)/m}$ since $k>\epsilon m$.
This completes the uniform imbalance estimate, including full
occupancy.

\subsection{Summing cycle moments over scales}
Apply \eqref{n:ideal-recursion} at root with \(G=1\) and \(l\) distinguished cards. At each height \(j\ge1\) (nodes of size \(2^j=2m\)), distinguished counts \(k_v\) at the nodes sum to \(l\), so \eqref{n:color-imbalance} costs only \(e^{Cl}\) over the tree for fixed \(\theta\). It remains to obtain, in the true \(H_N\) runs, for some small absolute \(\theta>0\)
\begin{equation}
\mathbb E 2^{\theta\sum_v c_v}\le e^{C l}
\label{n:all-cycle-moment}
\end{equation}
where \(c_v\) is the cycle count \(c_*\) of \eqref{n:node-colors} for the placements there.

Indeed, the logarithm of the noncycle factor at height $j$ is at
most $Clj/2^j$, and $\sum_{j\ge1}j2^{-j}<\infty$.
It is the cycle factor, rather than this deterministic imbalance
factor, that requires the remaining multiscale argument.

Condition on all input layers, fixing the assignments and inputs throughout the tree. At a node, given \(\alpha,\beta\) the two child permutations on columns, \(c_v\) is the number of cycles of \(\beta^{-1}\alpha\) entirely contained in the set \(D_v\) of input columns with both entering positions marked (distinguished). To see this, in a cycle of the two matchings one uses both inputs of participating columns, one going through each child. Following this circuit, going through the output matching from the \(\alpha\)-card at one such input column \(i\), and then through the input matching, brings us to the \(\alpha\)-card for column \(\beta^{-1}\alpha(i)\). Conversely a cycle of the comparison contained in \(D_v\) traces such a circuit, irrespective of the outermost output switches there. In particular \(c_v\) with inputs fixed depends just on smaller height output layers.

For \(h=1,2,4,\ldots\) group node heights \(j\) in bands \(h\le j<2h\). For such \(j\) consider conditioning on all inputs as above and also all output layers at heights \(<h/2\). For large \(h\), \(\alpha\) in the above comparison has still its last \(L=h/2\) switch layers (distinct axes in this child) fresh. Indeed these output layers of a height \(j-1\) child use output coins at heights \(j-L,\ldots,j-1\); each layer runs over all the \(m\) columns in that child. Even conditioning further on everything except these layers, \eqref{n:encoded-cycle-tail} applies to the comparison and \(D_v\), the latter of size at most \(k_v\). Consequently for sufficiently small \(\theta>0\),
\begin{equation}
\mathbb E\left[2^{2\theta h c_v}\mid {\rm inputs}, {\rm outputs}_{<h/2}\right]\le
\exp(C k_v h 2^{-\gamma' h})
\label{n:conditional-band-moment}
\end{equation}
for some absolute \(\gamma'>0\); in the conditioning we mean layer coins. Indeed for positive counts at or above \(k_v 2^{-\gamma L}\), multiplication by \(2^{2\theta h t}\) still gives geometric decay in \(ht\) from \eqref{n:encoded-cycle-tail}, and for counts below threshold one bounds the moment by \(2^{2\theta h k_v 2^{-\gamma L}}\). For \(k_v\ge1\) the additive exponentially small term fits as well by taking \(\gamma'\) small; \(k_v=0\) is automatic. For bounded \(h\), \eqref{n:conditional-band-moment} holds by enlarging the constant since \(c_v\le k_v\).

Conditional on the indicated layers, nodes on a fixed level have independent \(c_v\), using only their disjoint subtrees. So \eqref{n:conditional-band-moment} also bounds the \(2\theta h\) moment in this notation for the level sum, with \(k_v\) replaced by \(l\). H\"older over the at most \(h\) levels in the band gives
\[
\mathbb E\left[2^{2\theta\sum_{j\text{ in band}}\sum_{v\text{ at }j}c_v}\mid
{\rm inputs}, {\rm outputs}_{<h/2}\right]\le \exp(C l h 2^{-\gamma' h}).
\]
The next smaller band of the same parity of $\log_2h$ has heights
$h/4\le j<h/2$.  Its cycles depend only on output layers strictly
below $h/2$, so its exponential factor is measurable under the
conditioning for the band starting at $h$.  The tower property
therefore permits integrating each parity from the largest band
downwards.  It bounds that parity's $2\theta$ moment by
\[
 \exp\left\{Cl\sum_{h=1,2,4,\ldots}h2^{-\gamma'h}\right\}
 \le e^{C'l}.
\]
Cauchy--Schwarz combines the two parities at exponent $\theta$ and
proves~\eqref{n:all-cycle-moment}.  Independence was used within a
height, while the tower property handles the dependence across
heights.

\begin{proof}[Completion of Theorem~\ref{n:strong-density}]
Inserting \eqref{n:node-colors} and \eqref{n:color-imbalance} into
\eqref{n:ideal-recursion}, then applying \eqref{n:all-cycle-moment},
gives \eqref{n:strong-injection-density} with $p=1+\theta$.
Decrease $\theta$ if necessary so that $p\le2$.  The proof applies
to either order at the root and throughout its descendants, so it
proves both orientations with the same absolute constants.

To change normalization, put $s_l=(N)_l/N^l$.  Then
\[
 N^{-l}\sum_y(N^lH_N(x,y))^p
 =s_l^{-(p-1)}\frac1{(N)_l}
       \sum_y((N)_lH_N(x,y))^p.
\]
The decreasing list $\log(1-i/N)$, $0\le i<N$, has mean
$N^{-1}\log(N!/N^N)\ge-1$, since $N!\ge(N/e)^N$.
The mean of its first $l$ terms is no smaller, so $s_l\ge e^{-l}$.
This proves \eqref{n:strong-density-power-normalization} after an
increase of $C$.  Both positive products are self-adjoint, hence
$H_N(x,y)=H_N(y,x)$; the column statements follow.  The cases
$l=0$ and $N=1$ are immediate.
\end{proof}

\subsection{A fixed number of compositions gives a pointwise bound}

\begin{corollary}[Strong pointwise smoothing]
\label{n:strong-smoothing}
There are an absolute integer $u\ge1$ and an absolute $C<\infty$
such that, for every $N=2^d$, every coordinate order, every
$0\le l\le N$, every $x,y\in\mathcal X_{N,l}$, and either
$H_N=T_N^*T_N$ or $H_N=T_NT_N^*$,
\begin{equation}\label{n:strong-smoothing-eq}
 H_N^u(x,y)\le\frac{e^{Cl}}{(N)_l}.
\end{equation}
In particular, with unnormalized regular trace on $S_N$,
\begin{equation}\label{n:strong-regular-smoothing}
 \operatorname{Tr}_{\mathrm{reg}}|T_N|^{2u}\le e^{CN}.
\end{equation}
\end{corollary}

\begin{proof}
All function norms in this proof use the probability measure
$u_{N,l}$.  Write $A=e^{C_0l}$ for a common bound on the $L^p$
norms of the transition densities supplied by
Theorem~\ref{n:strong-density}.  The kernel formula and Minkowski's
inequality give $\|H_N\|_{1\to p}\le A$.  Self-adjointness gives
$\|H_N\|_{p'\to\infty}\le A$, where $p'=p/(p-1)$.
As a Markov operator, $H_N$ has $\infty$-to-$\infty$ norm one.
Riesz--Thorin interpolation consequently gives, for every $b\ge1$,
\[
 \|H_N\|_{b\to pb}\le A^{1/b}.
\]
Let $r$ be the least nonnegative integer with $p^r\ge p'$.
Composing the bounds for $b=1,p,\ldots,p^{r-1}$ yields
\[
 \|H_N^r\|_{1\to p^r}
 \le A^{\sum_{j=0}^{r-1}p^{-j}}.
\]
Since the underlying measure is a probability measure,
$\|f\|_{p'}\le\|f\|_{p^r}$.  One further application of
$H_N:p'\to\infty$ therefore proves
$\|H_N^{r+1}\|_{1\to\infty}\le e^{Cl}$.
Applying this to $(N)_l\mathbf1_{\{x\}}$, whose $L^1$ norm is
one, gives \eqref{n:strong-smoothing-eq} with $u=r+1$.

For $l=N$, the injection module is the regular permutation module
on $S_N$.  Every diagonal entry of $(T_N^*T_N)^u$ is at most
$e^{CN}/N!$, so its unnormalized trace is at most $e^{CN}$.
Since $|T_N|^{2u}=(T_N^*T_N)^u$, this is
\eqref{n:strong-regular-smoothing}.
\end{proof}

\subsection{The lower-diagram dimension factor}

The density theorem also distinguishes permutations of the tracked
labels.  We now make that distinction explicit.  This is useful
when a bound involving only the number of boxes below the first row
does not yet control the dimension of the lower diagram.

\begin{proposition}[Strong tail-dimension bound]
\label{n:strong-tail}
There are absolute constants $a>0$ and $C<\infty$ such that, for
every dyadic $N$, every coordinate order, and every partition
$\lambda=(N-k,\beta)\vdash N$, with
$\beta=(\lambda_2,\lambda_3,\ldots)\vdash k$,
\begin{equation}\label{n:strong-tail-eq}
 \|T_N(\lambda)\|_{\mathrm{op}}^2
       \le e^{Ck}D_\beta^{-a}.
\end{equation}
Here $D_\beta=\dim V_\beta$, including $D_\varnothing=1$.
One may take $a=1/p'=(p-1)/p$ for the exponent in
Theorem~\ref{n:strong-density}.
\end{proposition}

\begin{proof}
The case $k=0$ is immediate.  On ordered $k$-tuples, permutations
of the $k$ coordinate labels act on the right and commute with all
position permutations.  Choose an ordering $x_A$ of each unordered
$k$-set $A$.  The fiber over $A$ is then the regular representation
of $S_k$.  Since $H_N$ commutes with this right action, its block
from one fiber to another acts by convolution.  In the Fourier
decomposition of one fiber, the $\beta$ isotype is
$V_\beta\otimes V_\beta^*$, with the right action on the second
factor.  The block acts as $M_{AB}\otimes I$; it acts on the first,
multiplicity factor and may mix the copies of the right type.
Up to inverse or transpose conventions that do not affect the norm,
its multiplicity matrix is
\[
 M_{AB}=\sum_{\pi\in S_k}
                  H_N(x_A,x_B\pi)\rho_\beta(\pi).
\]
For unit vectors $v,w\in V_\beta$, Schur orthogonality
\cite[Proposition~3.14(ii)]{etingof} and unitarity give
\[
 \frac1{k!}\sum_{\pi\in S_k}
       |\langle v,\rho_\beta(\pi)w\rangle|^2=D_\beta^{-1},
 \qquad |\langle v,\rho_\beta(\pi)w\rangle|\le1.
\]
Because $p'\ge2$, the normalized $L^{p'}$ norm of this matrix
coefficient is at most $D_\beta^{-1/p'}$.  H\"older therefore gives
\begin{equation}\label{n:strong-slot-block}
 \|M_{AB}\|_{\mathrm{op}}
 \le k!N^{-k}D_\beta^{-1/p'}
 \left\{\frac1{k!}\sum_{\pi\in S_k}
       (N^kH_N(x_A,x_B\pi))^p\right\}^{1/p}.
\end{equation}

To sum this estimate, put $w_0=k!N^{-k}$ and denote the braces
to the power $1/p$ by $F_{AB}$.  There are $\binom Nk$ fibers,
so $\sum_Bw_0=(N)_k/N^k\le1$.  The density estimate gives
\[
 \sum_Bw_0 F_{AB}^p
 =N^{-k}\sum_{y\in\mathcal X_{N,k}}
             (N^kH_N(x_A,y))^p\le e^{Ck}.
\]
Weighted H\"older shows that every row sum of the scalar matrix
$(\|M_{AB}\|_{\mathrm{op}})$ is at most
$e^{Ck}D_\beta^{-1/p'}$.  The same bound holds for column sums:
self-adjointness and relabeling equivariance give
\[
 H_N(x_A,x_B\pi)=H_N(x_B,x_A\pi^{-1}).
\]
The Schur test now bounds the full operator on the right-$\beta$
isotypic space by $e^{Ck}D_\beta^{-1/p'}$.  Explicitly, if
$b_{AB}=\|M_{AB}\|_{\mathrm{op}}$ and both scalar sums are at most
$L$, then
\[
 \sum_A\Big\|\sum_BM_{AB}v_B\Big\|^2
 \le\sum_A\Big(\sum_Bb_{AB}\Big)
                    \sum_Bb_{AB}\|v_B\|^2
 \le L^2\sum_B\|v_B\|^2.
\]

As a left $S_N$ module, this right-$\beta$ isotypic space consists
of $D_\beta$ copies of
\[
 \operatorname{Ind}_{S_{N-k}\times S_k}^{S_N}
                    (\mathbf1\otimes V_\beta).
\]
The skew diagram $\lambda/\beta$ is a horizontal strip of size
$N-k$: its interlacing conditions are
$\lambda_i\ge\beta_i\ge\lambda_{i+1}$, and here
$\beta_i=\lambda_{i+1}$.  The horizontal Pieri rule therefore
places $V_\lambda$ in this induced module.  The preceding operator
bound applies to its $\lambda$ block.  Taking
$H_N=T_N^*T_N$ and using
$\|H_N(\lambda)\|_{\mathrm{op}}=
\|T_N(\lambda)\|_{\mathrm{op}}^2$ proves the assertion.
The other orientation gives the same singular-value conclusion.
\end{proof}
